\section{Decisiones y supuestos abiertos}\label{sec:pendientes}

Las decisiones 1 a 6, 12 a 15, 19 y 20 quedaron cerradas por el equipo. Las demás son tareas o puntos por confirmar. Los parámetros del modelo siguen siendo provisionales hasta validarlos con un actor de gestión del riesgo.

\begin{longtable}{@{}p{0.9cm}p{8cm}p{4.2cm}@{}}
\toprule
\textbf{N.º} & \textbf{Asunto} & \textbf{Estado} \\
\midrule
\endhead
1 & Fórmula de confianza (RM-09 y RM-10): coeficientes y valores de $F$ y $\tau$. & Decidido: se mantiene. Validar con un experto. \\
2 & Pesos de las dimensiones (RM-06 y RM-07). & Decidido: un perfil general fijo, ajustable. Validar con un experto. \\
3 & Tratamiento de la accesibilidad. & Decidido: alerta, sin penalizar la prioridad. \\
4 & Cifras de rendimiento (RNF-06). & Decidido: no se fijan cifras. \\
5 & Reportes anónimos. & Decidido: no existen. Todo bajo usuarios registrados (RF-04). \\
6 & Costo del despliegue en Railway (RNF-09). & Resuelto: el equipo ya tiene una cuenta de pago. \\
7 & Datos de barrios y viviendas del Atlántico (RD-02). & En recopilación. \\
8 & Revisión jurídica de la Ley 1581 de 2012 para datos personales y sensibles (RD-14 a RD-20). & Pendiente antes de cualquier uso real. \\
9 & Validación con una entidad de gestión del riesgo. & Sin definir. \\
10 & Equipo, mentora, institución y fecha de la portada. & Resuelto. \\
11 & Referencias: completar las tres sin DOI y leer los artículos citados antes de usarlos como argumento. & Pendiente. \\
12 & Identificación del solicitante: nombre, teléfono y sector obligatorios; documento opcional y obligatorio para entregas de mayor valor; vulnerabilidad con descripción de la situación. & Decidido (RF-02, RF-13). \\
13 & Quién entrega: brigadistas y organizaciones registradas. & Decidido. \\
14 & La entrega debe validarse sin conexión. & Decidido (RF-24, RNF-14). \\
15 & Tecnología del código: se propone un código de un solo uso con huella con sal y un registro encadenado por hash, en lugar de una cadena de bloques propiamente dicha (sección~7). & Decidido: se confirma, sin cadena de bloques real. \\
16 & Qué es una entrega ``de mayor valor'' para exigir el documento (RF-02). & Por definir. \\
17 & Vigencia del código de entrega y política de conservación de los datos (RD-20). & Por definir. \\
18 & Fuentes concretas de datos oficiales de eventos y daños (RD-04). & Por identificar y confirmar. \\
19 & Alcance de las 12 horas: se recortó el agrupamiento y los escenarios ``qué pasaría si'' a \Could{}. & Decidido; el resto de los \Must{} se hace en versión mínima. \\
20 & Alcance de la solución: registrar donaciones, priorizar necesidades y optimizar los recursos humanitarios. & Decidido (sección~1). \\
21 & Donaciones monetarias: solo se registran; el sistema no procesa pagos (RF-60). & Propuesta; por confirmar. \\
22 & Si el donante puede ocultar su nombre en las vistas públicas (RD-23). & Propuesta; por confirmar. \\
23 & Parámetros del plan de asignación: peso $w_j$, $\kappa$, $\lambda$ y equidad $\theta$ (sección~\ref{sec:optimizacion}). & Provisionales; validar con un experto. \\
24 & Quién gestiona los centros de acopio y cómo se registran (organizaciones aprobadas o la entidad coordinadora). & Por definir. \\
\bottomrule
\end{longtable}
\addtocounter{table}{-1}
